\documentclass[12pt, letterpaper]{article}

\usepackage[utf8]{inputenc}
\usepackage[T1]{fontenc}
\usepackage[spanish]{babel}
\usepackage{graphicx}
\usepackage{booktabs}

\usepackage[top=2.5cm, bottom=2.5cm, left=3cm, right=3cm]{geometry}

\usepackage[hidelinks]{hyperref}

\usepackage{tocloft}

\usepackage{listings}
\usepackage{xcolor}

\lstset{
    basicstyle=\ttfamily\small,
    keywordstyle=\color{blue},
    commentstyle=\color{gray},
    stringstyle=\color{orange},
    breaklines=true,
    frame=single,
    language=C
}

\usepackage{setspace}
\onehalfspacing %espaciado 1.5

\begin{document}

\begin{titlepage}
    \centering
    \vspace*{2cm}

    {\Large \textbf{Tecnológico de Costa Rica}}\\[0.5cm]
    {\large Principios de Sistemas Operativos [IC-6600]}\\[0.3cm]
    {\large Sede: Cartago}\\[0.3cm]
    {\large Profesor: Armando Arce Orozco}\\[2cm]

    {\LARGE \textbf{Proyecto 3: Sistema de archivos S3}}\\[2cm]

    {\large \textbf{Integrantes:}}\\[0.3cm]
    {\large Fernando Andrés Sánchez Hidalgo \hspace{0.5cm} [2022218688]}\\[0.2cm]
    {\large Kevin Espinoza Solano \hspace{0.5cm} [2023055841]}\\[2cm]

    {\large I Semestre 2026}\\[0.5cm]

    {\large \textbf{Repositorio:}}\\[0.2cm]
    \href{https://github.com/fash2110/OS-3_2026}{\texttt{https://github.com/fash2110/OS-3\_2026}}

    \vfill
\end{titlepage}

\tableofcontents
\newpage

\section{Introducción}
El proyecto tiene como objetivo simular un sistema de archivos basado en buckets similar al S3 (Simple Storage Service) de AWS. Para esto se implementan dos programas, un servidor que maneja operaciones de buckets y sus objetos (Creación, Modificación, Obtención y Eliminación) y recibe comandos del programa cliente. El programa cliente obtiene comandos del usuario por medio de la consola y los traslada para que el servidor realice las operaciones del sistema de archivos.

\section{Descripción del problema}
Para la implementación del proyecto, se generan dos programas escritos en C
\begin{itemize}
    \item Servidor: Enviar/recibir mensajes al/del cliente y realizar operaciones en sistema de archivos
    \item Cliente: Interfaz del usuario, envíar y recibir mensajes al servidor
\end{itemize}
La estrategia del programa servidor:
\begin{enumerate}
    \item Iniciar el servidor TCP en el puerto 8080 y quedar a la espera 
          de conexiones entrantes mediante \texttt{accept()}.
    \item Al conectarse un cliente, recibir la cabecera fija del mensaje 
          (\texttt{MsgHeader}) que indica el tipo de comando, flags y 
          longitudes de los campos.
    \item Leer secuencialmente del socket: el nombre del bucket, la clave 
          del objeto (\textit{key}), una segunda clave opcional (\textit{key2}) 
          y los datos binarios adjuntos, si los hay.
    \item Determinar el comando solicitado y despacharlo al manejador 
          correspondiente:
    \begin{itemize}
        \item \texttt{CMD\_LS}: listar buckets o el contenido de un bucket 
              filtrando por prefijo.
        \item \texttt{CMD\_MB}: crear un nuevo archivo \texttt{.bkt} que 
              representa el bucket.
        \item \texttt{CMD\_PUT}: almacenar un objeto en el bucket, 
              reutilizando espacio libre si existe o escribiendo al final 
              del archivo.
        \item \texttt{CMD\_GET}: recuperar el contenido de un objeto del 
              bucket y enviarlo al cliente.
        \item \texttt{CMD\_MV}: copiar el objeto al destino y eliminar el 
              original.
        \item \texttt{CMD\_RM}: eliminar uno o varios objetos, opcionalmente 
              de forma recursiva por prefijo.
        \item \texttt{CMD\_SYNC}: confirmar disponibilidad para recibir 
              múltiples objetos vía \texttt{CMD\_PUT}.
        \item \texttt{CMD\_RB}: eliminar un bucket, con verificación de 
              que esté vacío a menos que se use el flag \texttt{--force}.
    \end{itemize}
    \item Enviar al cliente una respuesta con cabecera \texttt{MsgHeader}, 
          indicando \texttt{RESP\_OK}, \texttt{RESP\_ERROR}, 
          \texttt{RESP\_DATA} o \texttt{RESP\_LIST} según corresponda, 
          junto con datos adjuntos si la operación lo requiere.
    \item Regresar al paso 2 para atender el siguiente comando del mismo 
          cliente, hasta que este cierre la conexión.
\end{enumerate}
Estrategia del programa cliente:
\begin{enumerate}
    \item Recibir el comando del usuario por línea de argumentos 
          (\texttt{argc}, \texttt{argv}) e identificar la operación 
          solicitada: \texttt{ls}, \texttt{mb}, \texttt{cp}, \texttt{mv}, 
          \texttt{rm}, \texttt{sync} o \texttt{rb}.
    \item Analizar los argumentos restantes, parsear las URIs con formato 
          \texttt{s3://bucket/key} para extraer el nombre del bucket y la 
          clave del objeto, e identificar las banderas opcionales 
          (\texttt{--recursive}, \texttt{--delete}, \texttt{--force}).
    \item Establecer una conexión TCP con el servidor en 
          \texttt{127.0.0.1:8080} mediante \\ \texttt{connect\_to\_server()}.
    \item Construir y enviar la cabecera \texttt{msgHeader} seguida del 
          nombre del bucket, la clave del objeto, una segunda clave 
          opcional (\textit{key2}) y los datos binarios adjuntos si la 
          operación lo requiere (por ejemplo, \texttt{cp} local $\to$ S3).
    \item Recibir la respuesta del servidor mediante 
          \texttt{recv\_response()}, que retorna el tipo de respuesta 
          (\texttt{RESP\_OK}, \texttt{RESP\_ERROR}, \texttt{RESP\_DATA}, 
          \texttt{RESP\_LIST}) junto con los datos adjuntos si los hay.
    \item Procesar la respuesta: mostrar el resultado al usuario, escribir 
          archivos locales en caso de descarga, o reportar el error 
          correspondiente.
    \item Cerrar la conexión con el servidor. Para operaciones que 
          requieren múltiples intercambios (como \texttt{sync} o 
          \texttt{cp --recursive}), se establece una nueva conexión por 
          cada operación individual.
\end{enumerate}

\section{Definición de estructuras de datos}
\subsection{Structs}
\subsubsection{Servidor}
\begin{lstlisting}
typedef struct __attribute__((packed)) {
    uint8_t  cmd;          /* tipo de comando / respuesta  */
    uint8_t  flags;        /* bits auxiliares (--recursive=1, --force=2, --delete=4) */
    uint16_t bucket_len;   /* longitud del nombre del bucket (sin '\0') */
    uint16_t key_len;      /* longitud de la clave/prefijo   (sin '\0') */
    uint16_t key2_len;     /* clave destino en mv            (sin '\0') */
    uint64_t data_len;     /* longitud de los datos adjuntos */
} msgHeader;
\end{lstlisting}
Cabecera que precede a todos los mensajes recibidos o emitidos por el servidor, contiene información como el tipo de comando enviado, banderas auxiliares, longitud del identificador del bucket, longitud del nombre de los objetos incluidos y longitud de los datos. \\

\begin{lstlisting}
typedef struct {
    char     magic[4];     // firma "BKT1" para archivos binarios
    uint32_t max_entries;
    uint32_t num_entries;  // entradas de directorio ocupadas
    uint32_t max_free;
    uint32_t num_free;     // entradas de free-list ocupadas
    uint64_t data_start;   // offset donde empieza la zona de datos
} bucketHeader;
\end{lstlisting}
Cabecera de cada bucket para gestión de sus objetos correspondientes. Todo bucket empieza con la firma \texttt{BKT1} por ser un archivo binario. En la cabecera se mantiene información como la capacidad máxima del bucket, cantidad de entradas ocupadas, capacidad máxima de espacios libres, cantidad de espacios libres actuales y un offset del inicio del espacio de datos calculado con \texttt{data\_start=sizeof(BucketHeader) + MAX\_ENTRIES \times sizeof(DirEntry) + MAX\_FREE \times sizeof(FreeBlock)} para garantizar que los metadatos y contenido del bucket no contengan intersección. \\

\begin{lstlisting}
typedef struct {
    char     key[MAX_KEY_LEN];      // clave/prefijo completo, "" si está libre
    uint64_t offset;                // byte donde empieza el contenido
    uint64_t size;                  // tamaño en bytes del objeto
    int      used;                  // 1 = slot ocupado, 0 = libre
} dirEntry;
\end{lstlisting}
Entrada de directorio de un objeto en el bucket para búsqueda directa del contenido, se tiene un campo para el identificador del objeto (ej: \texttt{fotos/vacaciones/2024.jpg} o "" si es un espacio libre), un offset que representa la distancia desde el inicio del archivo hasta el primer byte del contenido para búsqueda con \texttt{fseek()}, el tamaño del contenido del objeto y un bit para definir si es un campo libre u ocpudao. \\

\begin{lstlisting}
typedef struct {
    uint64_t offset;
    uint64_t size;
    int      used;   // 1 = hueco libre válido, 0 = slot de la lista sin usar
} freeBlock;
\end{lstlisting}
Espacio libre en el contenido del bucket, contiene el offset para su búsqueda directa en la zona de datos, tamaño del espacio libre para identificar si es lo suficientemente grande para el contenido que lo va a reemplazar y un bit que define si es un espacio válido u ocupado.

\subsubsection{Cliente}
\begin{lstlisting}
typedef struct __attribute__((packed)) {
    uint8_t  cmd;
    uint8_t  flags;
    uint16_t bucket_len;
    uint16_t key_len;
    uint16_t key2_len;
    uint64_t data_len;
} msgHeader;
\end{lstlisting}
La estructura \texttt{msgHeader} del cliente es idéntica a la del servidor 
y debe mantenerse así por diseño: ambos extremos de la comunicación deben 
interpretar los mismos bytes de la misma forma. El cliente la utiliza 
exclusivamente para construir y enviar mensajes, mientras que el servidor 
la utiliza para recibirlos e interpretarlos.

\subsection{Variables globales}
Para ambos programas de cliente y servidor, se utilizan las mismas variables globales, ya que se requiere la misma configuración para una comunicación exitosa.
\begin{itemize}
    \item MAX\_KEY\_LEN: Tamaño máximo para nombre de un object
    \item MAX\_ENTRIES: Capacidad máxima de un bucket
    \item MAX\_FREE: Capacidad máxima del espacio libre
    \item BUCKETS\_DIR: Directorio donde se almacenan los buckets del lado del servidor (\texttt{./buckets})
    \item SERVER\_PORT: Puerto para recibir mensajes (\texttt{8080})
    \item BACKLOG: Cantidad de clientes que permite el servidor para operaciones secuenciales (\textbf{10})
    \item enum de comandos de mensajes
    \begin{itemize}
        \item CMD\_LS: 1
        \item CMD\_MB: 2
        \item CMD\_PUT: 3
        \item CMD\_GET: 4
        \item CMD\_MV: 5
        \item CMD\_RM: 6
        \item CMD\_SYNC: 7
        \item CMD\_RB: 8
    \end{itemize}
    \item Códigos de respuesta
    \begin{itemize}
        \item RESP\_OK: 100
        \item RESP\_ERROR: 101
        \item RESP\_DATA: 102
        \item RESP\_LIST: 103
    \end{itemize}
    \item Banderas auxiliares
    \begin{itemize}
        \item FLAG\_RECURSIVE: 0x01
        \item FLAG\_FORCE: 0x02
        \item FLAG\_DELETE: 0x04
    \end{itemize}
\end{itemize}

\section{Componentes principales}

\subsection*{Servidor}
\begin{itemize}
    \item \texttt{main(void) $\to$ int}
        \begin{itemize}
            \item Inicializa el socket TCP en el puerto 8080
            \item Queda en espera de conexiones entrantes y delega cada 
                  cliente a \texttt{handle\_client()}
        \end{itemize}
    \item \texttt{handle\_client(int client\_fd) $\to$ void}
        \begin{itemize}
            \item Componente central del servidor
            \item Recibe la cabecera \texttt{MsgHeader} y los campos del 
                  mensaje (bucket, key, key2, datos)
            \item Despacha el comando al manejador correspondiente según 
                  el campo \texttt{cmd}
            \item Mantiene el ciclo de atención hasta que el cliente 
                  cierra la conexión
        \end{itemize}

    \subsection*{Manejadores de comandos}
    \item \texttt{handle\_ls(int fd, char* bucket, char* prefix) $\to$ void}
        \begin{itemize}
            \item Lista todos los buckets disponibles si \texttt{bucket} 
                  está vacío
            \item Lista los objetos de un bucket filtrando por prefijo
        \end{itemize}
    \item \texttt{handle\_mb(int fd, char* bucket) $\to$ void}
        \begin{itemize}
            \item Crea un nuevo archivo \texttt{.bkt} con su cabecera, 
                  directorio y lista de espacios libres inicializados
        \end{itemize}
    \item \texttt{handle\_put(int fd, char* bucket, char* key, 
          unsigned char* data, uint64\_t data\_len) $\to$ void}
        \begin{itemize}
            \item Almacena un objeto en el bucket
            \item Reutiliza un \texttt{FreeBlock} disponible mediante 
                  \textit{first-fit}, o escribe al final del archivo
            \item Crea el bucket automáticamente si no existe
        \end{itemize}
    \item \texttt{handle\_get(int fd, char* bucket, char* key) $\to$ void}
        \begin{itemize}
            \item Recupera el contenido de un objeto del bucket y lo 
                  envía al cliente como \texttt{RESP\_DATA}
        \end{itemize}
    \item \texttt{handle\_mv(int fd, char* bucket\_src, char* key\_src, 
          char* bucket\_dst, char* key\_dst) $\to$ void}
        \begin{itemize}
            \item Copia el objeto al bucket destino y elimina el original
            \item Crea el bucket destino automáticamente si no existe
        \end{itemize}
    \item \texttt{handle\_rm(int fd, char* bucket, char* key, 
          int recursive) $\to$ void}
        \begin{itemize}
            \item Elimina un objeto específico o, en modo recursivo, 
                  todos los objetos cuya clave inicie con el prefijo dado
        \end{itemize}
    \item \texttt{handle\_rb(int fd, char* bucket, int force) $\to$ void}
        \begin{itemize}
            \item Elimina un bucket verificando que esté vacío
            \item Con \texttt{force=1} elimina el bucket aunque contenga 
                  objetos
        \end{itemize}

    \subsection*{Gestión de buckets}
    \item \texttt{bucket\_create(char* bucket\_name) $\to$ int}
        \begin{itemize}
            \item Crea el archivo \texttt{.bkt} e inicializa el 
                  \texttt{BucketHeader}, el arreglo de \texttt{DirEntry} 
                  y la lista de \texttt{FreeBlock}
        \end{itemize}
    \item \texttt{bucket\_put\_object(char* bucket, char* key, 
          unsigned char* data, uint64\_t size) $\to$ int}
        \begin{itemize}
            \item Busca si el objeto ya existe para reemplazarlo
            \item Aplica el algoritmo \textit{first-fit} sobre la lista 
                  de \texttt{FreeBlock} para asignar espacio
            \item Actualiza la \texttt{DirEntry} correspondiente y la 
                  lista de espacios libres
        \end{itemize}
    \item \texttt{bucket\_get\_object(char* bucket, char* key, 
          uint64\_t* out\_size) $\to$ unsigned char*}
        \begin{itemize}
            \item Localiza el objeto por su clave en el directorio
            \item Retorna un buffer con el contenido; el llamador es 
                  responsable de liberar la memoria con \texttt{free()}
        \end{itemize}
    \item \texttt{bucket\_delete\_object(char* bucket, char* key) $\to$ int}
        \begin{itemize}
            \item Marca la \texttt{DirEntry} como libre (\texttt{used=0})
            \item Registra el espacio liberado como un nuevo 
                  \texttt{FreeBlock}
        \end{itemize}
    \item \texttt{bucket\_remove(char* bucket, int force) $\to$ int}
        \begin{itemize}
            \item Verifica si el bucket contiene objetos antes de 
                  eliminarlo
            \item Elimina el archivo \texttt{.bkt} del sistema de archivos
        \end{itemize}

    \subsection*{Helpers de red}
    \item \texttt{net\_send\_all(int fd, void* buf, size\_t len) $\to$ int}
        \begin{itemize}
            \item Garantiza el envío completo de \texttt{len} bytes al 
                  socket, reintentando si \texttt{send()} transmite menos 
                  bytes de los esperados
        \end{itemize}
    \item \texttt{net\_recv\_all(int fd, void* buf, size\_t len) $\to$ int}
        \begin{itemize}
            \item Garantiza la recepción completa de \texttt{len} bytes 
                  del socket, reintentando si \texttt{recv()} entrega 
                  menos bytes de los esperados
        \end{itemize}
    \item \texttt{send\_response(int fd, uint8\_t type, char* msg) $\to$ void}
        \begin{itemize}
            \item Envía una respuesta simple al cliente con cabecera 
                  \texttt{MsgHeader} y un mensaje de texto opcional
        \end{itemize}
    \item \texttt{send\_data\_response(int fd, uint8\_t type, 
          void* data, uint64\_t len) $\to$ void}
        \begin{itemize}
            \item Envía una respuesta con datos binarios adjuntos, 
                  utilizada en operaciones \texttt{GET} y \texttt{LS}
        \end{itemize}
\end{itemize}

\subsection*{Cliente}
\begin{itemize}
    \item \texttt{main(int argc, char* argv[]) $\to$ int}
        \begin{itemize}
            \item Identifica el comando solicitado comparando 
                  \texttt{argv[1]} contra los comandos disponibles
            \item Desplaza \texttt{argv} para que cada subcomando reciba 
                  sus argumentos desde \texttt{argv[0]}
            \item Delega la ejecución al comando correspondiente e imprime 
                  el uso correcto si el comando es desconocido
        \end{itemize}

    \subsection*{Comandos}
    \item \texttt{cmd\_ls(int argc, char* argv[]) $\to$ int}
        \begin{itemize}
            \item Si se provee una URI \texttt{s3://bucket[/prefijo]}, 
                  lista los objetos del bucket filtrando por prefijo
            \item Sin argumentos, solicita al servidor la lista de todos 
                  los buckets disponibles
        \end{itemize}
    \item \texttt{cmd\_mb(int argc, char* argv[]) $\to$ int}
        \begin{itemize}
            \item Parsea la URI \texttt{s3://bucket} y envía 
                  \texttt{CMD\_MB} al servidor para crear el bucket
        \end{itemize}
    \item \texttt{cmd\_cp(int argc, char* argv[]) $\to$ int}
        \begin{itemize}
            \item Determina la dirección de la copia según si los 
                  argumentos son URIs S3 o rutas locales:
            \begin{itemize}
                \item Local $\to$ S3: lee el archivo con 
                      \texttt{read\_file()} y lo envía con 
                      \texttt{do\_put()}. Con \texttt{--recursive} 
                      recorre el directorio con \texttt{opendir()} y 
                      sube cada archivo manteniendo la ruta relativa 
                      como clave
                \item S3 $\to$ local: descarga el objeto con 
                      \texttt{do\_get()}. Con \texttt{--recursive} 
                      obtiene primero el listado del bucket y descarga 
                      cada objeto individualmente
                \item S3 $\to$ S3: descarga el objeto origen con 
                      \texttt{CMD\_GET} y lo sube al destino con 
                      \texttt{CMD\_PUT} en una nueva conexión
            \end{itemize}
        \end{itemize}
    \item \texttt{cmd\_mv(int argc, char* argv[]) $\to$ int}
        \begin{itemize}
            \item Local $\to$ S3: reutiliza la lógica de \texttt{cmd\_cp} 
                  y elimina el archivo local con \texttt{remove()}
            \item S3 $\to$ local: descarga el objeto con \texttt{do\_get()} 
                  y lo elimina del bucket con \texttt{CMD\_RM}
            \item S3 $\to$ S3: envía \texttt{CMD\_MV} con 
                  \texttt{key2 = "bucket\_dst\textbackslash nkey\_dst"} 
                  para que el servidor realice la operación atómicamente
        \end{itemize}
    \item \texttt{cmd\_rm(int argc, char* argv[]) $\to$ int}
        \begin{itemize}
            \item Envía \texttt{CMD\_RM} con la bandera 
                  \texttt{FLAG\_RECURSIVE} activa si se especifica 
                  \texttt{--recursive}, eliminando todos los objetos 
                  cuya clave inicie con el prefijo dado
        \end{itemize}
    \item \texttt{cmd\_sync(int argc, char* argv[]) $\to$ int}
        \begin{itemize}
            \item Local $\to$ S3: recorre el directorio local y sube 
                  cada archivo con \texttt{do\_put()}. Con 
                  \texttt{--delete} obtiene el listado del bucket y 
                  elimina los objetos remotos que ya no existen localmente
            \item S3 $\to$ local: obtiene el listado del bucket y 
                  descarga únicamente los objetos cuyo tamaño difiere 
                  del archivo local correspondiente, omitiendo los que 
                  ya están sincronizados
        \end{itemize}
    \item \texttt{cmd\_rb(int argc, char* argv[]) $\to$ int}
        \begin{itemize}
            \item Envía \texttt{CMD\_RB} con la bandera 
                  \texttt{FLAG\_FORCE} activa si se especifica 
                  \texttt{--force}, permitiendo eliminar buckets que 
                  aún contengan objetos
        \end{itemize}

    \subsection*{Helpers de red y sistema de archivos}
    \item \texttt{connect\_to\_server(void) $\to$ int}
        \begin{itemize}
            \item Crea un socket TCP y establece la conexión con el 
                  servidor en \texttt{SERVER\_HOST:SERVER\_PORT}
            \item Retorna el file descriptor del socket o \texttt{-1} 
                  en caso de error
        \end{itemize}
    \item \texttt{send\_header(int fd, uint8\_t cmd, uint8\_t flags, 
          char* bucket, char* key, char* key2, uint64\_t data\_len) $\to$ int}
        \begin{itemize}
            \item Construye y envía la cabecera \texttt{msgHeader} 
                  seguida de los campos \texttt{bucket}, \texttt{key} 
                  y \texttt{key2}, calculando automáticamente las 
                  longitudes de cada campo
        \end{itemize}
    \item \texttt{recv\_response(int fd, unsigned char** out\_data, 
          uint64\_t* out\_len) $\to$ int}
        \begin{itemize}
            \item Lee la cabecera de respuesta del servidor y, si 
                  \texttt{data\_len > 0}, recibe los datos adjuntos en 
                  un buffer asignado dinámicamente
            \item El llamador es responsable de liberar el buffer con 
                  \texttt{free()}
        \end{itemize}
    \item \texttt{do\_put(int fd, char* local, char* bucket, 
          char* key) $\to$ int}
        \begin{itemize}
            \item Lee el archivo local con \texttt{read\_file()}, envía 
                  la cabecera \texttt{CMD\_PUT} seguida del contenido 
                  binario y verifica la respuesta del servidor
        \end{itemize}
    \item \texttt{do\_get(int fd, char* bucket, char* key, 
          char* local) $\to$ int}
        \begin{itemize}
            \item Envía \texttt{CMD\_GET} y escribe el contenido 
                  recibido en la ruta local con \texttt{write\_file()}
            \item Si \texttt{local} es un directorio, construye el 
                  nombre del archivo usando la parte final de la clave
        \end{itemize}
    \item \texttt{parse\_s3\_uri(char* uri, char* bucket, 
          char* key) $\to$ int}
        \begin{itemize}
            \item Divide una URI \texttt{s3://bucket/key} en sus 
                  componentes \texttt{bucket} y \texttt{key}
            \item Retorna \texttt{0} si la URI es válida o \texttt{-1} 
                  si no tiene el prefijo \texttt{s3://}
        \end{itemize}
    \item \texttt{read\_file(char* path, uint64\_t* out\_size) 
          $\to$ unsigned char*}
        \begin{itemize}
            \item Lee el contenido completo de un archivo local en un 
                  buffer dinámico; el llamador debe liberar la memoria 
                  con \texttt{free()}
        \end{itemize}
    \item \texttt{write\_file(char* path, unsigned char* data, 
          uint64\_t size) $\to$ int}
        \begin{itemize}
            \item Escribe datos en una ruta local, creando los 
                  directorios intermedios necesarios con \texttt{mkdir()}
        \end{itemize}
\end{itemize}

\section{Mecanismo de creación de archivos y comunicación con el servidor}

\subsection{Estructura del archivo \texttt{.bkt}}
Cada bucket se representa como un único archivo binario con extensión 
\texttt{.bkt} almacenado en el directorio \texttt{./buckets/} del servidor. 
Al crear un bucket con \texttt{bucket\_create()}, se escribe una estructura 
fija compuesta por tres zonas contiguas:

\begin{enumerate}
    \item \textbf{Cabecera} (\texttt{bucketHeader}): ocupa los primeros 
          bytes del archivo e incluye la firma \texttt{BKT1}, las 
          capacidades máximas y el offset \texttt{data\_start} que delimita 
          el inicio de la zona de datos.
    \item \textbf{Directorio de objetos}: arreglo de \texttt{MAX\_ENTRIES} 
          entradas \texttt{dirEntry}, cada una inicializada con 
          \texttt{used=0}. Ocupa \texttt{MAX\_ENTRIES $\times$ 
          sizeof(dirEntry)} bytes inmediatamente después de la cabecera.
    \item \textbf{Lista de espacios libres}: arreglo de \texttt{MAX\_FREE} 
          entradas \texttt{freeBlock}, cada una inicializada con 
          \texttt{used=0}. Ocupa \texttt{MAX\_FREE $\times$ 
          sizeof(freeBlock)} bytes a continuación del directorio.
    \item \textbf{Zona de datos}: región donde se escribe el contenido 
          binario de los objetos, a partir del offset \texttt{data\_start}.
\end{enumerate}

\subsection{Algoritmo de asignación de espacio (first-fit)}
Cuando se almacena un objeto nuevo o se reemplaza uno existente con un 
tamaño diferente, la función \texttt{allocate\_space()} recorre la lista 
de \texttt{freeBlock} en orden y selecciona el primer hueco cuyo campo 
\texttt{size} sea mayor o igual al tamaño requerido. Si no existe ningún 
hueco suficiente, el objeto se escribe al final del archivo, calculando 
el offset con \texttt{data\_end\_offset()}. Una vez utilizado un 
\texttt{freeBlock}, su campo \texttt{used} se marca como \texttt{0} para 
liberarlo de la lista. Cuando un objeto es eliminado o reemplazado, su 
espacio se registra como un nuevo \texttt{freeBlock} mediante 
\texttt{add\_free\_block()}.

\subsection{Protocolo de comunicación cliente-servidor}
La comunicación se realiza sobre una conexión TCP entre el cliente y el 
servidor en el puerto \texttt{8080}. Todo intercambio de mensajes sigue 
el mismo formato:

\begin{enumerate}
    \item El cliente envía una cabecera fija de 20 bytes 
          (\texttt{msgHeader}) con el tipo de comando, banderas y las 
          longitudes de los campos que siguen.
    \item A continuación envía secuencialmente: el nombre del bucket, la 
          clave del objeto (\textit{key}), una segunda clave opcional 
          (\textit{key2}, usada en \texttt{CMD\_MV}), y los datos binarios 
          adjuntos si \texttt{data\_len > 0}.
    \item El servidor procesa el comando y responde con otra 
          \texttt{msgHeader} indicando el resultado 
          (\texttt{RESP\_OK}, \texttt{RESP\_ERROR}, \texttt{RESP\_DATA} 
          o \texttt{RESP\_LIST}), seguida de los datos adjuntos si los hay.
\end{enumerate}

Las funciones \texttt{net\_send\_all()} y \texttt{net\_recv\_all()} 
garantizan la transmisión completa de cada bloque de bytes, reintentando 
la llamada al sistema si \texttt{send()} o \texttt{recv()} transfieren 
menos bytes de los solicitados, lo cual es habitual en sockets TCP. Esto 
permite transmitir correctamente tanto texto como archivos binarios 
(imágenes, ejecutables, etc.).

\section{Pruebas}

\subsection{Metodología}
Las pruebas fueron realizadas desde un sistema windows utilizando la herramienta de WSL (Windows subsystem for Linux) utilizando la siguiente configuración en la figura \ref{fig:wsl_config}. Primeramente se levanta el servidor utilizando
\begin{lstlisting}
    >> cd servidor
    >> ./aws-s3_server
\end{lstlisting}
Una vez se tiene el servidor levantado, se pueden empezar a realizar las pruebas desde el directorio cliente
\begin{lstlisting}
    >> cd ..
    >> cd cliente
    >> cd pruebas
    >> ./benchmark.sh
\end{lstlisting}

\begin{figure}[h]
    \centering
    \includegraphics[width=1\linewidth]{WSLconfig.png}
    \caption{Configuración de la maquina virtual WSL}
    \label{fig:wsl_config}
\end{figure}

Se generó un script de pruebas que puede ser encontrado en el repositorio de github en la ubicación \texttt{/cliente/pruebas/benchmark.sh}, dentro del script de pruebas se encuentran las siguientes pruebas de rendimiento en la tabla \ref{tab:pruebas}

\begin{table}[h]
\caption{Pruebas de rendimiento}
\label{tab:pruebas}
\centering
\begin{tabular}{p{3cm}p{10cm}}
\toprule
\textbf{Prueba} & \textbf{Descripción} \\
\midrule
\texttt{latencia}        & crear, listar (con y sin bucket) y eliminar bucket \\ \\
\texttt{subida}          & subida de archivos (10KB, 512KB, 5MB, 20MB) \\ \\
\texttt{descarga}        & descarga de archivos (S3 $\to$ local) \\ \\
\texttt{sync}            & sincronización (local $\to$ S3) \\ \\
\texttt{move + remove}   & mover archivos dentro del bucket, eliminar objeto en el bucket, eliminar objetos recursivamente \\ \\
\texttt{sobreescritura}  & sobreescribir archivos con mismo y diferente tamaño \\ \\
\bottomrule
\end{tabular}
\end{table}

\subsection{Resultados}
Una vez ejecutado el script de pruebas, se genera un archivo \texttt{/cliente/pruebas/} \texttt{resultados\_benchmark.txt} en el cual se documenta el resultado de las pruebas realizadas, los resultados se pueden visualizar en la tabla \ref{tab:resultados} 

\begin{table}[h]
\caption{Resultados de pruebas}
\label{tab:resultados}
\centering
\begin{tabular}{p{10cm}p{3cm}}
\toprule
\textbf{Prueba} & \textbf{tiempo (ms)} \\
\midrule
\texttt{latencia}            & \\
\texttt{mb}                  & 2 \\
\texttt{ls (listar buckets)} & 3 \\
\texttt{ls (listar objetos)} & 2 \\
\texttt{rb (vacío)}          & 3 \\
\midrule
\texttt{subida}             & \\
\texttt{10 KB}              & 3 \\
\texttt{512KB}              & 2 \\
\texttt{5MB}                & 3 \\
\texttt{20MB}               & 3 \\
\midrule
\texttt{descarga}           & \\
\texttt{10 KB}              & 3 \\
\texttt{512KB}              & 2 \\
\texttt{5MB}                & 3 \\
\texttt{20MB}               & 2 \\
\midrule
\texttt{sync}               & \\
\texttt{10 archivos}        & 2 \\
\texttt{50 archivos}        & 2 \\
\texttt{100 archivos}       & 3 \\
\midrule
\texttt{move + remove}      & \\
\texttt{mv (mismo bucket)}  & 2 \\
\texttt{rm (objeto)}        & 3 \\
\texttt{rm --recursive}     & 2 \\
\midrule
\texttt{sobreescritura}     & \\
\texttt{mismo tamaño}       & 2 \\
\texttt{distinto tamaño}    & 2 \\
\bottomrule
\end{tabular}
\end{table}

Se puede visualizar en todas las pruebas que los tiempo de respuesta se centran entre 2 y 3ms, esto se debe en gran parte a que ambos programas cliente y servidor son ejectuados dentro del mismo sistema. En general, los resultados demuestran que la arquitectura de archivo único por bucket con directorio de entradas de tamaño fijo permite tiempos de acceso predecibles y constantes, sin cambios observables al aumentar el tamaño de los objetos, o la cantidad de archivos sincronizados en el entorno de prueba utilizado.

\section{Conclusiones}
\begin{itemize}
    \item La representación de cada bucket como un único archivo binario 
          \texttt{.bkt} con una zona de metadatos fija al inicio y una 
          zona de datos variable al final demostró ser una estrategia 
          eficaz para simular un sistema de almacenamiento de objetos, 
          evitando la dependencia del sistema de archivos del sistema 
          operativo para la organización interna de los objetos. Las 
          pruebas de latencia base confirmaron esto, con operaciones de 
          metadatos como \texttt{mb}, \texttt{ls} y \texttt{rb} 
          completando consistentemente en 2--3 ms.

    \item El uso del atributo \texttt{\_\_attribute\_\_((packed))} en la 
          estructura \texttt{msgHeader} fue esencial para garantizar que 
          la cabecera ocupe exactamente 20 bytes en ambos extremos de la 
          comunicación, eliminando problemas de alineación de memoria que 
          podrían corromper el protocolo entre cliente y servidor.

    \item El algoritmo \textit{first-fit} implementado en 
          \texttt{allocate\_space()} permite reutilizar los espacios 
          liberados por objetos eliminados o reemplazados, evitando el 
          crecimiento indefinido del archivo \texttt{.bkt}. Las pruebas 
          de sobreescritura confirmaron que tanto la sustitución de objetos 
          del mismo tamaño como de distinto tamaño se realizan en 2 ms, 
          sin diferencia medible entre ambos casos en entorno local. Sin 
          embargo, el no dividir un bloque libre mayor en dos partes tras 
          una asignación parcial puede generar fragmentación interna con 
          el tiempo.

    \item Los tiempos de subida y descarga se mantuvieron estables entre 
          2 y 3 ms independientemente del tamaño del archivo transferido, 
          desde 10 KB hasta 20 MB. Dado que las pruebas se ejecutaron con 
          cliente y servidor en la misma máquina sobre el adaptador de 
          loopback (\texttt{127.0.0.1}), el factor dominante fue el acceso 
          a disco y no la latencia de red, lo que indica que el rendimiento 
          en un entorno de red real podría variar significativamente.

    \item Las pruebas de \texttt{sync} mostraron tiempos consistentes de 
          2--3 ms al escalar de 10 a 100 archivos, lo que sugiere que el 
          costo por conexión TCP en entorno local es despreciable. Sin 
          embargo, dado que \texttt{cmd\_sync()} establece una conexión 
          independiente por cada archivo, en una red real con latencia 
          apreciable este comportamiento podría convertirse en un cuello 
          de botella, siendo una mejora natural el uso de conexiones 
          persistentes o el envío en grupos.

    \item La separación clara entre el programa cliente y el servidor 
          permite que el cliente sea una interfaz liviana que únicamente 
          construye mensajes y presenta resultados, mientras que toda la 
          lógica de gestión de archivos reside en el servidor, lo que 
          facilita la extensión futura del sistema.

    \item El servidor atiende clientes de forma secuencial, por lo que 
          operaciones como \texttt{sync} o \texttt{cp --recursive}, que 
          requieren múltiples conexiones, son procesadas una a la vez. 
          Una mejora natural sería introducir concurrencia mediante hilos 
          o procesos hijos para atender múltiples clientes simultáneamente, 
          especialmente considerando que las pruebas se realizaron en un 
          entorno de un único cliente y los tiempos podrían degradarse 
          considerablemente bajo carga concurrente real.
\end{itemize}

\end{document}